\exposetitle{III}{Infinitesimal extensions}
\label{III}

\begin{center}
\textit{by M.~Demazure}
\end{center}

{\renewcommand{\thefootnote}{(0)}\nde{Version of 14/10/2024}}%
In this expos\'e, one places oneself in the following general situation.
One has a scheme~\(S\) and a coherent nilpotent ideal~\(\mathcal{I}\)
on~\(S\). One denotes by~\(S_n\) the closed subscheme of~\(S\) defined
by the ideal~\(\mathcal{I}^{n+1}\) (\(n\geqslant 0\)). In particular
\(S_0\) is defined by~\(\mathcal{I}\). As \(\mathcal{I}\) is nilpotent,
\(S_n\) equals~\(S\) for \(n\) large enough, and the~\(S_i\) have the same
underlying topological space. A typical example of this situation is the
following: \(S\) is the spectrum of a local artinian ring~\(A\),
\(\mathcal{I}\) is the ideal defined by the radical of~\(A\), hence
\(S_0\) is the spectrum of the residue field of~\(A\).

In the preceding situation, one is given a certain number of data
above~\(S_0\), and one seeks above~\(S\) data which lift them, that is to
say, which give them back by base change from~\(S\) to~\(S_0\). This is
done step by step, by means of the~\(S_n\). At each step, one proposes
to define the obstructions encountered and to classify, when they exist,
the solutions obtained.

The passage from~\(S_n\) to~\(S_{n+1}\) may be generalized as follows:
one has a scheme~\(S\), two ideals \(\mathcal{I}\) and~\(\mathcal{J}\)
with \(\mathcal{I}\supset\mathcal{J}\), and
\(\mathcal{I}\cdot\mathcal{J}=0\) (in the preceding case, \(S\),
\(\mathcal{I}\) and~\(\mathcal{J}\) are respectively \(S_{n+1}\),
\(\mathcal{I}/\mathcal{I}^{n+2}\),
\(\mathcal{I}^{n+1}/\mathcal{I}^{n+2}\)). One denotes by~\(S_0\)
(\resp \(S_{\mathcal{J}}\)) the closed subscheme of~\(S\) defined
by~\(\mathcal{I}\) (\resp \(\mathcal{J}\)), and one poses a problem of
extension from~\(S_{\mathcal{J}}\) to~\(S\).

In SGA~1~III were treated the problems of extension of morphisms of
schemes and of extension of schemes. We shall pose here the problems of
extension of morphisms of groups, of extension of group structures and
of extension of subgroups.

We have gathered in an n\textsuperscript{o}~0 the results of SGA~1~III
which will be useful to us, so as to put them in the form most practical
for our purpose, and so as to spare the reader from having constantly to
refer to SGA~1~III.%
{\renewcommand{\thefootnote}{(1)}\nde{The reader may prefer to begin by
reading section~1, which is easier, and which can serve as motivation
and as a guide for the results obtained in section~0.}}
The n\textsuperscript{o}~1 gathers computations of cohomology of groups,
useful in the sequel, and which have nothing to do with the theory of
schemes. Numbers~2 and~3 treat respectively of \emph{the extension of
morphisms of groups and of the extension of group structures}. In
n\textsuperscript{o}~4, we have briefly recalled the proof of a result
stated in TDTE~IV concerning the extension of subschemes, and applied
this result to the problem of extension of subgroups. \emph{For the
sequel of the Seminar, only the result of n\textsuperscript{o}~2,
concerning the extension of morphisms of groups, will be
indispensable.}%
{\renewcommand{\thefootnote}{(2)}\nde{However, 3.10 is used in XXIV,
1.13. On the other hand, 4.37 is used in the proof of IX, 3.2~bis
and 3.6~bis, but these also follow from results of Exp.~X, which do not
use~4.37.}}

The idea of reducing the problems of infinitesimal extensions to the
usual cohomology computations in group extensions was suggested by
J.~Giraud during the oral expos\'e (whose computations were distinctly
more complicated and less transparent). Unfortunately, it seems that
this method applies well only to the first two problems studied, and we
have been unable to avoid rather painful computations in the case of
extensions of subgroups.

To simplify the language, we shall call \(Y\)-functor, \resp
\(Y\)-scheme, etc., a functor, \resp scheme, etc., equipped with a
morphism into the functor~\(Y\), thus extending the definitions of
Expos\'e~I (which concerned only the case of a representable~\(Y\)).

\sgasection{0}{Reminders from SGA 1 III and various remarks}
\label{III.sec.0}

Let us first state a general definition.

\begin{definition}{0.1}
\label{III.0.1}
Let \(\mathcal{C}\) be a category, \(X\) an object of
\(\widehat{\mathcal{C}}\), \(G\) a \(\widehat{\mathcal{C}}\)-group
operating on~\(X\). One says that \(X\) is formally principal
homogeneous%
{\renewcommand{\thefootnote}{(3)}\nde{One also says ``pseudo-torsor'',
\cf EGA~\(\mathrm{IV}_{4}\), 16.5.15. On the other hand, the more
general notion of formally homogeneous object (not necessarily
principal homogeneous) is defined in Exp.~IV, 6.7.1.}}
under~\(G\) if the following equivalent conditions are satisfied:
\begin{enumeratei}
\item for each object \(S\) of~\(\mathcal{C}\), the set \(X(S)\) is empty
or principal homogeneous under \(G(S)\);
\item the morphism of functors \(G\times X\to X\times X\) defined
set-theoretically by \((g,x)\mapsto(gx,x)\) is an isomorphism.
\end{enumeratei}
\end{definition}

This done, we are going to put the results of SGA~1~III, \S~5%
{\renewcommand{\thefootnote}{(4)}\nde{See also
EGA~\(\mathrm{IV}_{4}\), 16.5.14--18.}}
in the form which will be the most useful to us. We shall employ the
following general notations throughout this number. One has a
scheme~\(S\) and on~\(S\) two quasi-coherent ideals \(\mathcal{I}\)
and~\(\mathcal{J}\) such that
\[
\mathcal{I}\supset\mathcal{J}
\qquad\text{and}\qquad
\mathcal{I}\cdot\mathcal{J}=0.
\]
One will therefore have in particular \(\mathcal{J}^{2}=0\). One will
denote by~\(S_0\) (\resp \(S_{\mathcal{J}}\)) the closed subscheme
of~\(S\) defined by the ideal~\(\mathcal{I}\)
(\resp \(\mathcal{J}\)). For every \(S\)-functor~\(X\), one will
systematically denote by \(X_0\) and~\(X_{\mathcal{J}}\) the functors
obtained by base change from~\(S\) to \(S_0\) and~\(S_{\mathcal{J}}\).
The same notations for a morphism.

\begin{definition}{0.1.1}
\label{III.0.1.1}
{\renewcommand{\thefootnote}{(5)}\nde{One has added the numbering
0.1.1, \ldots, 0.1.13, so as to bring out the definitions and results
which are found there.}}%
Let \(X\) be an \(S\)-functor. Let us define a functor~\(X^{+}\)
above~\(S\) by the formula
\[
\Hom_S(Y,X^{+})
=\Hom_{S_{\mathcal{J}}}(Y_{\mathcal{J}},X_{\mathcal{J}})
=\Hom_S(Y_{\mathcal{J}},X)
\]
for a variable \(S\)-scheme~\(Y\). In the notations of Exp.~II, 1, one
has
\[
X^{+}\simeq\prod_{S_{\mathcal{J}}/S}X_{\mathcal{J}}.
\]
The identity morphism of~\(X_{\mathcal{J}}\) defines by construction an
\(S\)-morphism
\[
p_X\colon X\longrightarrow X^{+}.
\]
{\renewcommand{\thefootnote}{(6)}\nde{One has added the following
sentence, and spelled out the two remarks which follow.}}%
Explicitly, for every \(S\)-scheme~\(Y\), the map
\[
p_X(Y)\colon
\Hom_S(Y,X)\longrightarrow\Hom_S(Y,X^{+})=\Hom_S(Y_{\mathcal{J}},X)
\]
is the map induced by the morphism \(Y_{\mathcal{J}}\to Y\).
\end{definition}

\begin{remark}{0.1.2}
\label{III.0.1.2}
Let us now remark that if \(X\) is an \(S\)-functor in groups,
\(X_{\mathcal{J}}\) is an \(S_{\mathcal{J}}\)-functor in groups; then
the defining formula of~\(X^{+}\) equips it with a structure of
\(S\)-functor in groups, and the description of~\(p_X\) above shows that
\(p_X\colon X\to X^{+}\) is a morphism of \(S\)-functors in groups.
\end{remark}

\begin{remark}{0.1.3}
\label{III.0.1.3}
On the other hand, for every \(S\)-functor in groups~\(Y\), one has
\[
\Hom_{S\text{-gr.}}(Y,X^{+})
=\Hom_{S_{\mathcal{J}}\text{-gr.}}(Y_{\mathcal{J}},X_{\mathcal{J}}).
\]
Indeed, let \(f\in\Hom_S(Y,X^{+})\), corresponding to
\(f_{\mathcal{J}}\in\Hom_{S_{\mathcal{J}}}(Y_{\mathcal{J}},X_{\mathcal{J}})\);
the condition for \(f\in\Hom_{S\text{-gr.}}(Y,X^{+})\) is that, for every
\(T\to S\) and \(y,y'\in Y(T)\), one has
\(f(y\cdot y')=f(y)\cdot f(y')\), and this is equivalent to
\[
f_{\mathcal{J}}(y_{\mathcal{J}})\cdot f_{\mathcal{J}}(y'_{\mathcal{J}})
=f_{\mathcal{J}}((y\cdot y')_{\mathcal{J}});
\]
as \((y\cdot y')_{\mathcal{J}}=y_{\mathcal{J}}\cdot y'_{\mathcal{J}}\)
(since \(Y(T)\to Y(T_{\mathcal{J}})=Y_{\mathcal{J}}(T_{\mathcal{J}})\) is
a morphism of groups), this is the condition for \(f_{\mathcal{J}}\) to
be a morphism of groups. Applying this to \(Y=X\), one finds again that
\(p_X\), which corresponds to \(\mathrm{id}_{X_{\mathcal{J}}}\), is a
morphism of \(S\)-functors in groups.
\end{remark}

Let us now return to the general case, but suppose that \(X\) is an
\(S\)-scheme. An \(S\)-morphism of a variable \(S\)-scheme~\(Y\)
into~\(X^{+}\) being by definition an \(S_{\mathcal{J}}\)-morphism
\(g_{\mathcal{J}}\) of~\(Y_{\mathcal{J}}\) into~\(X_{\mathcal{J}}\), one
is going to define an \(X^{+}\)-functor in abelian groups~\(L_X\) as
follows.

\begin{scholium}{0.1.4}
\label{III.0.1.4}
{\renewcommand{\thefootnote}{(7)}\nde{One has added this scholium.}}%
If \(\pi\colon Y\to S\) is a morphism of schemes and \(\mathcal{M}\) an
\(\OS\)-module, one denotes by \(\mathcal{M}\otimes_{\OS}\OY\) the
inverse image \(\pi^{*}(\mathcal{M})\). If \(\mathcal{J}\) is an ideal
of~\(\OS\), one denotes by \(\mathcal{J}\OY\) the ideal of~\(\OY\), image
of the morphism
\[
\pi^{*}(\mathcal{J})=\mathcal{J}\otimes_{\OS}\OY\longrightarrow\OY.
\]
\end{scholium}

Let us note that, for every morphism of \(S\)-schemes
\(f\colon Z\to Y\), one has an epimorphism of \(\mathcal{O}_Z\)-modules
\[
(0.1.4)\qquad
f^{*}(\mathcal{J}\OY)\longrightarrow\mathcal{J}\mathcal{O}_Z,
\]
as follows from the commutative diagram below:
\[
\xymatrix{
\mathcal{J}\otimes_{\OS}\OY\otimes_{\OY}\mathcal{O}_Z \ar[d]
  & \mathcal{J}\otimes_{\OS}\mathcal{O}_Z \ar[d] \\
f^{*}(\mathcal{J}\OY)=(\mathcal{J}\OY)\otimes_{\OY}\mathcal{O}_Z \ar[r]
  & \mathcal{J}\mathcal{O}_Z.
}
\]

\begin{definition}{0.1.5}
\label{III.0.1.5}
{\renewcommand{\thefootnote}{(8)}\nde{The introduction of the functor
\(L_X\) leads to an extension ``functorial in \(Y\)'' (and in \(X\)) of
SGA~1, III, Prop.~5.1; see below 0.1.8 and~0.1.9, as well as 0.1.10
and~0.2. Moreover, when \(X\) is an \(S\)-functor in groups, \(L_X\)
gives rise, under certain hypotheses, to an \(S\)-functor in groups
\(L'_X\) and to an exact sequence \(1\to L'_X\to X\to X^{+}\) (\cf 0.4,
0.9 and~0.11), which play an essential role in this expos\'e (\cf
Theorems~2.1, 3.5 and~4.21).}}%
Let \(X\) be an \(S\)-scheme. For every \(X^{+}\)-scheme~\(Y\), given by a
morphism \(g_{\mathcal{J}}\colon Y_{\mathcal{J}}\to X_{\mathcal{J}}\),
one sets
\[
\Hom_{X^{+}}(Y,L_X)
=\Hom_{\mathcal{O}_{Y_0}}\bigl(g_0^{*}(\Omega^1_{X_0/S_0}),\mathcal{J}\OY\bigr),
\]
where \(\Omega^1_{X_0/S_0}\) denotes the module of relative differentials
of~\(X_0\) with respect to~\(S_0\) (\cf SGA~1, I.1 or
EGA~\(\mathrm{IV}_{4}\), 16.3), and where one regards \(\mathcal{J}\OY\)
as an \(\mathcal{O}_{Y_0}\)-module thanks to the fact that it is
annihilated by~\(\mathcal{I}\).
\end{definition}

Then, \(L_X\) is an \(X^{+}\)-functor in abelian groups.%
{\renewcommand{\thefootnote}{(9)}\nde{One has added what follows.}}
Indeed, for every \(X^{+}\)-morphism \(f\colon Z\to Y\), the
functor~\(f_0^{*}\) and the morphism
\(f_0^{*}(\mathcal{J}\OY)\to\mathcal{J}\mathcal{O}_Z\) of~(0.1.4) induce a
natural morphism of abelian groups \(L_X(f)\):
\begin{align*}
\Hom_{\mathcal{O}_{Y_0}}\bigl(g_0^{*}(\Omega^1_{X_0/S_0}),\mathcal{J}\OY\bigr)
&\longrightarrow
\Hom_{\mathcal{O}_{Z_0}}\bigl(f_0^{*}g_0^{*}(\Omega^1_{X_0/S_0}),
f_0^{*}(\mathcal{J}\OY)\bigr)\\
&\longrightarrow
\Hom_{\mathcal{O}_{Z_0}}\bigl(f_0^{*}g_0^{*}(\Omega^1_{X_0/S_0}),
\mathcal{J}\mathcal{O}_Z\bigr).
\end{align*}

Finally, let us remark that \(L_X(f)\) is described in a simple way
locally, as follows. Let us first note that \(Y\) and~\(Y_{\mathcal{J}}\)
have the same underlying topological space, and that the same holds for
\(V\) and~\(V_{\mathcal{J}}\), if \(V\) is an open of~\(Y\). Let then
\(U=\Spec(A)\) be an affine open of~\(X\) above an affine open
\(\Spec(\Lambda)\) of~\(S\), \(V=\Spec(B)\) an affine open of~\(Y\) such
that \(g_{\mathcal{J}}(V_{\mathcal{J}})\subset U\), and \(W=\Spec(C)\) an
affine open of~\(f^{-1}(V)\). Let \(J\) and~\(I\) be the ideals
of~\(\Lambda\) corresponding to \(\mathcal{J}\) and~\(\mathcal{I}\). Then
\(f\) (\resp \(g_{\mathcal{J}}\)) induces a morphism of
\(\Lambda\)-algebras \(\theta\colon B\to C\)
(\resp \(\phi\colon A\to B/JB\)), and one evidently has
\(\theta(JB)\subset JC\). On the other hand,
\(m\in\Hom_{\mathcal{O}_{Y_0}}\bigl(g_0^{*}(\Omega^1_{X/S}),\mathcal{J}\OY\bigr)\)
induces an element~\(D\) of
\begin{align*}
\Hom_{\mathcal{O}_{V_0}}\bigl(g_0^{*}(\Omega^1_{U/S}),\mathcal{J}\mathcal{O}_V\bigr)
&=\Hom_{B/IB}\bigl(\Omega^1_{A/\Lambda}\otimes_A B/IB,\,JB\bigr)\\
&=\operatorname{Der}_{\Lambda}(A,JB),
\end{align*}
and the image of \(L_X(f)(m)\) in
\begin{align*}
\Hom_{\mathcal{O}_{W_0}}\bigl(f_0^{*}g_0^{*}(\Omega^1_{U/S}),
\mathcal{J}\mathcal{O}_Z\bigr)
&=\Hom_{C/IC}\bigl(\Omega^1_{A/\Lambda}\otimes_A C/IC,\,JC\bigr)\\
&=\operatorname{Der}_{\Lambda}(A,JC)
\end{align*}
is none other than \(\theta\circ D\).
% typo?: printed \(\Omega^1_{X/S}\), \(\Omega^1_{U/S}\) (not
% \(\Omega^1_{X_0/S_0}\), \(\Omega^1_{U_0/S_0}\)) and
% \(\mathcal{J}\mathcal{O}_Z\) in the \(\mathcal{O}_{W_0}\)-Hom; kept as printed

\begin{remark}{0.1.6}
\label{III.0.1.6}
{\renewcommand{\thefootnote}{(10)}\nde{One has spelled out this
remark.}}%
Let \(f\colon X\to W\) be an \(S\)-morphism. It induces an \(S\)-morphism
\(f^{+}\colon X^{+}\to W^{+}\) defined as follows: if \(g\) is an element
of \(\Hom_S(Y,X^{+})\), corresponding to an \(S\)-morphism
\(g_{\mathcal{J}}\colon Y_{\mathcal{J}}\to X\), then \(f^{+}(g)\) is the
element \(f\circ g_{\mathcal{J}}\) of \(\Hom_S(Y,W^{+})\). It is clear
that the following diagram is commutative:
\[
\xymatrix{
X \ar[r]^{f} \ar[d]_{p_X} & W \ar[d]^{p_W} \\
X^{+} \ar[r]^{f^{+}} & W^{+}.
}
\]
\end{remark}

\begin{enoncerm}{Reminders}{0.1.7}
\label{III.0.1.7}
{\renewcommand{\thefootnote}{(11)}\nde{One has added this paragraph
(see also \textup{[DG70]}, \S~I.4, n\textsuperscript{os}~1--2).}}%
In this paragraph, given an \(S\)-scheme~\(X\), one ``recalls'' certain
functorial properties of the module of differentials
\(\Omega^1_{X/S}\) and of the first infinitesimal neighborhood of the
diagonal, \(\Delta^{(1)}_{X/S}\), properties which follow easily from
EGA~\(\mathrm{IV}_{4}\), \S\S~16.1--16.4, but which do not appear there
explicitly.

a)~Let us begin by recalling the following facts (\cf EGA~II,
\S\S~1.2--1.5). Let \(g\colon Y\to X\) be a morphism of schemes,
\(\pi\colon X'\to X\) an \emph{affine} \(X\)-scheme, \(\mathcal{B}\) the
quasi-coherent \(\OX\)-algebra \(\pi_{*}(\mathcal{O}_{X'})\); then the
\(Y\)-scheme \(Y\times_X X'\) is affine and corresponds to the
quasi-coherent \(\OY\)-algebra \(g^{*}(\mathcal{B})\), and one has a
commutative diagram of bijections:
\[
\xymatrix{
\Hom_X(Y,X') \ar[r]^{\sim} \ar[d]_{\wr}
  & \Hom_Y(Y,Y\times_X X') \ar[d]^{\wr} \\
\Hom_{\OX\text{-alg.}}(\mathcal{B},g_{*}(\OY)) \ar[r]^{\sim}
  & \Hom_{\OY\text{-alg.}}(g^{*}(\mathcal{B}),\OY).
}
\]
Moreover, these bijections are functorial in the pair \((X,X')\),
\ie, if \(W'\) is an affine \(W\)-scheme, corresponding to the
quasi-coherent \(\mathcal{O}_W\)-algebra~\(\mathcal{A}\), if one has a
commutative diagram of morphisms of schemes
\[
\xymatrix{
& X' \ar[r]^{f'} & W' \\
Y \ar[ur]^{g'} \ar[r]_{g} & X \ar[r]_{f} & W
}
\]
and if one denotes by \(\phi\colon\mathcal{A}\to f_{*}(\mathcal{B})\) and
\(\phi^{\sharp}\colon f^{*}(\mathcal{A})\to\mathcal{B}\)
(\resp \(\theta\colon\mathcal{B}\to g_{*}(\OY)\) and
\(\theta^{\sharp}\colon g^{*}(\mathcal{B})\to\OY\)) the algebra morphisms
associated with~\(f'\) (\resp with a variable \(X\)-morphism
\(g'\colon Y\to X'\)), then one has the following commutative diagram
(where \(Y\) is viewed as a \(W\)-scheme via \(f\circ g\)):
\[
\xymatrix@C=1.6em@R=2.6em{
\Hom_X(Y,X') \ar[r]^-{g'\mapsto f'\circ g'} \ar[dd]_{\wr}
  & \Hom_W(Y,W') \ar[dr]^{\sim}
  & \\
&
  & \Hom_{\mathcal{O}_W\text{-alg.}}(\mathcal{A},f_{*}g_{*}(\OY))
    \ar[dl]_{\wr} \\
\Hom_{\OX\text{-alg.}}(\mathcal{B},g_{*}(\OY))
  \ar[r]^-{\theta\mapsto\theta\circ\phi^{\sharp}}
  \ar[d]_{\wr}
  \ar[urr]^-{\theta\mapsto f_{*}(\theta)\circ\phi}
  & \Hom_{\OX\text{-alg.}}(f^{*}(\mathcal{A}),g_{*}(\OY)) \ar[d]^{\wr}
  & \\
\Hom_{\OY\text{-alg.}}(g^{*}(\mathcal{B}),\OY)
  \ar[r]^-{\theta^{\sharp}\mapsto\theta^{\sharp}\circ g^{*}(\phi^{\sharp})}
  & \Hom_{\OY\text{-alg.}}(g^{*}f^{*}(\mathcal{A}),\OY).
  &
}
\]

b)~Let now \(X\) be an \(S\)-scheme. Let \(\Omega^1_{X/S}\) be the module
of differentials of~\(X\) over~\(S\), and \(\Delta^{(1)}_{X/S}\) the first
infinitesimal neighborhood of the diagonal immersion
\(\delta_X\colon X\to X\times_S X\); it is a subscheme of
\(X\times_S X\), of which the diagonal \(\Delta_{X/S}\) is a closed
subscheme. One denotes by \(\mathrm{pr}^{i}_X\) (\(i=1,2\)) the two
projections of \(X\times_S X\), and by~\(\pi_X\) the restriction of
\(\mathrm{pr}^{1}_X\) to~\(\Delta^{(1)}_{X/S}\).

On the one hand, every morphism \(f\colon X\to W\) of \(S\)-schemes
induces an \(S\)-morphism
\(\Delta^{(1)}f\colon\Delta^{(1)}_{X/S}\to\Delta^{(1)}_{W/S}\) such that
the following diagram is commutative:
\[
\xymatrix{
X \ar[r]^{\delta_X} \ar[d]_{f}
  & \Delta^{(1)}_{X/S} \ar[r] \ar[d]_{\Delta^{(1)}f}
  & X\times_S X \ar[r]^{\mathrm{pr}^{i}_X} \ar[d]_{f\times f}
  & X \ar[d]^{f} \\
W \ar[r]_{\delta_W}
  & \Delta^{(1)}_{W/S} \ar[r]
  & W\times_S W \ar[r]_{\mathrm{pr}^{i}_W}
  & W.
}
\]
On the other hand, \(\Delta^{(1)}_{X/S}\) is, via the projection
\(\pi_X\), an affine \(X\)-scheme, spectrum of the augmented
quasi-coherent \(\OX\)-algebra
\[
\mathcal{P}^{1}_{X/S}=\OX\oplus\Omega^1_{X/S},
\]
where \(\Omega^1_{X/S}\) is an ideal of square zero; the augmentation is
the morphism of \(\OX\)-algebras
\(\varepsilon_X\colon\mathcal{P}^{1}_{X/S}\to\OX\) which vanishes
on~\(\Omega^1_{X/S}\) and which corresponds to the closed immersion
\(\delta_X\colon X\hookrightarrow\Delta^{(1)}_{X/S}\). Then, every
morphism of \(S\)-schemes \(f\colon X\to W\) induces a morphism of
augmented \(\OX\)-algebras
\[
f^{*}(\mathcal{P}^{1}_{W/S})
=\OX\oplus f^{*}(\Omega^1_{W/S})
\longrightarrow
\mathcal{P}^{1}_{X/S}
=\OX\oplus\Omega^1_{X/S},
\]
\ie, in an equivalent way, a morphism of \(\OX\)-modules
\[
f_{X/W/S}\colon f^{*}(\Omega^1_{W/S})\longrightarrow\Omega^1_{X/S},
\]
\cf EGA~\(\mathrm{IV}_{4}\) (16.4.3.6) (and (16.4.18.2) for the notation
\(f_{X/W/S}\)).

As \(\pi_X\colon\Delta^{(1)}_{X/S}\to X\) is affine, then, by~a),
\(\Delta^{(1)}f\) is entirely determined by~\(f_{X/W/S}\) and, for every
\(X\)-scheme \(g\colon Y\to X\), the set
\[
\Hom_X\bigl(Y,\Delta^{(1)}_{X/S}\bigr)
\simeq
\Hom_{\OY\text{-alg.}}\bigl(\OY\oplus g^{*}(\Omega^1_{X/S}),\OY\bigr)
\]
is identified with a subset of
\(\Hom_{\OY}\bigl(g^{*}(\Omega^1_{X/S}),\OY\bigr)\), namely the subset
\[
\Hom^{\square}_{\OY}\bigl(g^{*}(\Omega^1_{X/S}),\OY\bigr)
\]
formed of the \(\OY\)-morphisms
\(\psi\colon g^{*}(\Omega^1_{X/S})\to\OY\) such that \(\Im(\psi)\) is an ideal
of~\(\OY\) of square zero.%
{\renewcommand{\thefootnote}{(12)}\nde{In the proof of SGA~1, III~5.1,
one must accordingly correct the sentence ``Now the algebra
homomorphisms \(\mathcal{A}\to\OY\) correspond~\ldots'' (the sequel of
\loccit being unchanged).}}

Consequently, applying~a) to the diagram
\[
\xymatrix{
& \Delta^{(1)}_{X/S} \ar[r]^{\Delta^{(1)}f} & \Delta^{(1)}_{W/S} \\
Y \ar[ur]^{g'} \ar[r]_{g} & X \ar[r]_{f} & W
}
\]
and taking into account the fact that \(\Delta^{(1)}f\) is the restriction
to~\(\Delta^{(1)}_{X/S}\) of \(f\times f\), one obtains the following
commutative diagram, functorial in the \(X\)-scheme
\(Y\xrightarrow{g}X\):
% typo?: top right target printed as \(\Hom_W(Y,\Delta^{(1)}_{W/S})\),
% the same object as the middle right target; kept as printed
\[
\xymatrix@C=2.2em@R=1.8em{
& \Hom_X(Y,X\times_S X)
  \ar[rr]^-{(g,g')\mapsto(f\circ g,\,f\circ g')}
  && \Hom_W\bigl(Y,\Delta^{(1)}_{W/S}\bigr) \\
\textup{\scriptsize(0.1.7 (*))}
  & \Hom_X\bigl(Y,\Delta^{(1)}_{X/S}\bigr)
  \ar[rr]^-{g'\mapsto(\Delta^{(1)}f)\circ g'}
  \ar@{^{(}->}[u]
  && \Hom_W\bigl(Y,\Delta^{(1)}_{W/S}\bigr)
  \ar@{^{(}->}[u] \\
& \Hom^{\square}_{\OY}\bigl(g^{*}(\Omega^1_{X/S}),\OY\bigr)
  \ar[rr]^-{\psi\mapsto\psi\circ g^{*}(f_{X/W/S})}
  \ar[u]_{\simeq}
  && \Hom^{\square}_{\OY}\bigl(g^{*}f^{*}(\Omega^1_{W/S}),\OY\bigr)
  \ar[u]^{\simeq}
}
\]
\end{enoncerm}

\begin{remark}{0.1.7.1}
\label{III.0.1.7.1}
Let us end this paragraph with the following remark, which will be useful
later (\cf 0.1.10). If one denotes by \(L^{\square}_X\) the \(X\)-functor
which, to every \(X\)-scheme \(g\colon Y\to X\), associates
\(\Hom^{\square}_{\OY}\bigl(g^{*}(\Omega^1_{X/S}),\OY\bigr)\), and by
\(L^{\square}_f\colon L^{\square}_X\to L^{\square}_W\) the morphism of
functors defined above (which, to every \(\psi\in L^{\square}_X(Y)\),
associates \(\psi\circ g^{*}(f_{X/W/S})\)), what precedes shows that one
has a commutative diagram of functors:
\[
\xymatrix{
X\times_X L^{\square}_X \ar[d]_{f\times L^{\square}_f}
  & \Delta^{(1)}_{X/S} \ar[l]_{\sim} \ar@{^{(}->}[r] \ar[d]_{\Delta^{(1)}f}
  & X\times_S X \ar[d]^{f\times f} \\
W\times_W L^{\square}_W
  & \Delta^{(1)}_{W/S} \ar[l]^{\sim} \ar@{^{(}->}[r]
  & W\times_S W.
}
\]
\end{remark}

\begin{theorem}{0.1.8}
\label{III.0.1.8}
(SGA~1, III~5.1)%
{\renewcommand{\thefootnote}{(13)}\nde{One has reproduced here the
statement and the proof of SGA~1, III~5.1, which one needs in order to
prove Proposition~0.2 below (see already Corollary~0.1.9 which follows).
On the other hand, \(\mathfrak{J}\) (\resp \(\mathfrak{I}\)) is another
calligraphy of the letter~J (\resp~I); returning to the earlier notations
(\(\mathcal{J}\subset\mathcal{I}\) ideals of \(\OS\) such that
\(\mathcal{I}\mathcal{J}=0\)), one will apply this to
\(\mathfrak{J}=\mathcal{J}\OY\)
(\resp \(\mathfrak{I}=\mathcal{I}\OY\)).}}
Let \(Y\), \(X\) be two \(S\)-schemes, \(\mathfrak{J}\) a quasi-coherent
ideal of~\(\OY\) of square zero, \(Y_{\mathfrak{J}}\) the closed subscheme
of~\(Y\) defined by~\(\mathfrak{J}\), and
\(g_{\mathfrak{J}}\colon Y_{\mathfrak{J}}\to X\) an \(S\)-morphism.

a)~The set \(\mathrm{P}(g_{\mathfrak{J}})\) of \(S\)-morphisms
\(g\colon Y\to X\) which extend~\(g_{\mathfrak{J}}\) is either empty, or
principal homogeneous under the abelian group
\[
\Hom_{\mathcal{O}_{Y_{\mathfrak{J}}}}
\bigl(g_{\mathfrak{J}}^{*}(\Omega^1_{X/S}),\mathfrak{J}\bigr).
\]

b)~If \(i\colon Y_0\hookrightarrow Y_{\mathfrak{J}}\) is the closed
immersion defined by a quasi-coherent ideal
\(\mathfrak{I}\supset\mathfrak{J}\) such that
\(\mathfrak{I}\mathfrak{J}=0\), and if
\(g_0=g_{\mathfrak{J}}\circ i\), the preceding abelian group is
isomorphic to
\[
\Hom_{\mathcal{O}_{Y_0}}
\bigl(g_0^{*}(\Omega^1_{X/S}),\mathfrak{J}\bigr).
\]
\end{theorem}

\begin{proof}
(b)~follows at once from~(a). Indeed, \(\mathfrak{J}\), being annihilated
by~\(\mathfrak{I}\), may be regarded as an \(\mathcal{O}_{Y_0}\)-module,
whence, by adjunction,
\[
\Hom_{\mathcal{O}_{Y_{\mathfrak{J}}}}
\bigl(g_{\mathfrak{J}}^{*}(\Omega^1_{X/S}),\mathfrak{J}\bigr)
=\Hom_{\mathcal{O}_{Y_0}}
\bigl(i^{*}g_{\mathfrak{J}}^{*}(\Omega^1_{X/S}),\mathfrak{J}\bigr).
\]
To prove~(a), one may suppose \(\mathrm{P}(g_{\mathfrak{J}})\neq\emptyset\),
\ie that there exists an \(S\)-morphism \(g\colon Y\to X\)
extending~\(g_{\mathfrak{J}}\). Let us denote by~\(j\) the immersion
\(Y_{\mathfrak{J}}\hookrightarrow Y\). Then \(\mathrm{P}(g_{\mathfrak{J}})\)
is the set of \(S\)-morphisms \(g'\colon Y\to X\) such that
\(g'\circ j=g_{\mathfrak{J}}\). The datum of such a~\(g'\) is equivalent
to the datum of an \(S\)-morphism
\[
h\colon Y\longrightarrow X\times_S X
\]
such that \(\mathrm{pr}_1\circ h=g\) and
\(h_{\mathfrak{J}}=\delta\circ g_{\mathfrak{J}}\), where
\(h_{\mathfrak{J}}=h\circ j\) and \(\delta\) is the diagonal immersion
\(X\hookrightarrow X\times_S X\):
\[
\xymatrix{
X\times_S X \ar[d]_{\mathrm{pr}_1}
  & Y_{\mathfrak{J}}
    \ar[l]_-{h_{\mathfrak{J}}=\delta\circ g_{\mathfrak{J}}}
    \ar[d]^{j} \\
X
  & Y \ar[l]_{g} \ar@{-->}[ul]_{h}.
}
\]
As \(h_{\mathfrak{J}}\) factors through~\(\delta\) and \(Y\) is in the first
infinitesimal neighborhood of the immersion
\(j\colon Y_{\mathfrak{J}}\to Y\) (since \(\mathfrak{J}^{2}=0\)), then, by
functoriality (\cf EGA~\(\mathrm{IV}_{4}\), 16.2.2~(i)), the sought-for
\(h\) factor, in a unique way, through~\(\Delta^{(1)}_{X/S}\)
(\cf 0.1.7). Let us set
\[
Y'=\Delta^{(1)}_{X/S}\times_X Y
\qquad\text{and}\qquad
Y'_{\mathfrak{J}}
=Y'\times_Y Y_{\mathfrak{J}}
=\Delta^{(1)}_{X/S}\times_X Y_{\mathfrak{J}}.
\]
Then the sought-for \(h\) are in bijection with the sections~\(u\) of
\(Y'\to Y\) which extend the section
\(u_{\mathfrak{J}}=(\delta\circ g_{\mathfrak{J}},\mathrm{id})\) of
\(Y'_{\mathfrak{J}}\to Y_{\mathfrak{J}}\). On the other hand, \(Y'\)
(\resp \(Y'_{\mathfrak{J}}\)) is a scheme \emph{affine over}~\(Y\)
(\resp \(Y_{\mathfrak{J}}\)), corresponding to the quasi-coherent algebra
\[
\mathcal{A}=\OY\oplus g^{*}(\Omega^1_{X/S}),
\qquad\text{resp.}\qquad
\mathcal{A}_{\mathfrak{J}}
=\mathcal{A}\otimes_{\OY}\mathcal{O}_{Y_{\mathfrak{J}}}
=\mathcal{O}_{Y_{\mathfrak{J}}}\oplus
g_{\mathfrak{J}}^{*}(\Omega^1_{X/S}).
\]
Let us denote by \(\varepsilon\colon\mathcal{A}\to\OY\) the canonical
augmentation of~\(\mathcal{A}\) (\ie the morphism of \(\OY\)-algebras
\(\mathcal{A}\to\OY\) which vanishes on \(g^{*}(\Omega^1_{X/S})\)), and let
us define likewise
\(\varepsilon_{\mathfrak{J}}\colon\mathcal{A}_{\mathfrak{J}}
\to\mathcal{O}_{Y_{\mathfrak{J}}}\). Then
\[
\Gamma(Y'/Y)\simeq\Hom_{\OY\text{-alg.}}(\mathcal{A},\OY),
\qquad
\Gamma(Y'_{\mathfrak{J}}/Y_{\mathfrak{J}})
\simeq\Hom_{\mathcal{O}_{Y_{\mathfrak{J}}}\text{-alg.}}
\bigl(\mathcal{A}_{\mathfrak{J}},\mathcal{O}_{Y_{\mathfrak{J}}}\bigr),
\]
and, via these isomorphisms, the section
\(u=(\delta\circ g,\mathrm{id})\) (\resp \(u_{\mathfrak{J}}\)) corresponds
to~\(\varepsilon\) (\resp \(\varepsilon_{\mathfrak{J}}\)).
\renewcommand{\qedsymbol}{}
\end{proof}
